\section{隐私权限模型：赛博分身应该如何被约束}

赛博分身要替用户主动面对世界，就必须有权限模型。一个分身可以浏览哪些 context，可以对谁说话，可以代表用户表达什么，可以自动完成哪些动作，哪些动作必须确认，这些都是产品底层问题。没有权限模型，分身越主动越危险。

\subsection{四层权限}

\noindent\begin{tabularx}{\textwidth}{p{0.22\textwidth}p{0.34\textwidth}X}
\toprule
权限层 & 分身能做什么 & 风险控制 \\
\midrule
读取层 & 读取用户资料、历史表达、偏好和关系 & 用户可查看、删除、分级授权。 \\
草稿层 & 生成评论、私信、推荐理由，但不发送 & 用户确认后发布。 \\
低风险行动层 & 点赞、收藏、公开低风险互动 & 频率限制、审计日志。 \\
高风险行动层 & 私聊、敏感观点、关系承诺、商业动作 & 必须人工确认，可撤回。 \\
\bottomrule
\end{tabularx}

\begin{knowledgebox}{权限模型也是产品体验}
用户愿意交出 context 的前提，是知道分身不会失控。清晰权限不是限制产品，而是让用户敢用产品。
\end{knowledgebox}

\subsection{Context 的所有权}

AI 社交最敏感的问题，是 context 到底属于谁。用户的聊天、作品、社交偏好和关系数据不能只是平台资产。长期看，用户需要可迁移、可导出、可删除、可本地保存的 context 权利。否则平台越懂用户，用户越不安。

\begin{importantbox}{Context 权利清单}
用户至少应拥有查看权、解释权、删除权、导出权、分级授权权和敏感动作确认权。AI 社交网络若想建立长期信任，必须把这些权利做成产品能力。
\end{importantbox}

\subsection{本章小结}

赛博分身越主动，权限模型越重要。AI 社交的信任基础，不是模型更会说话，而是用户能控制自己的 context 和分身行为。

